\begin{table}[htbp]
\centering
\caption{Empirical Application Results: No de-rounding, keeping coefficients reported as zero}
\label{tab:application_none_keepmu0}
\resizebox{\ifdim\width>\linewidth\linewidth\else\width\fi}{!}{% shrinks the table only if it is wider than the text
\begin{tabular}{lccccccccccccccccccc}
 &  &  &  &  & \multicolumn{2}{c}{$J_0=20$} & \multicolumn{2}{c}{$J_0=30$} & \multicolumn{7}{c}{EKW $p$-values} & \multicolumn{4}{c}{Kudrin (2025) $p$-values} \\
\cmidrule(lr){6-7} \cmidrule(lr){8-9} \cmidrule(lr){10-16} \cmidrule(lr){17-20}
 & $n$ & Articles & Med.\ $\nu$ & $\Delta_\nu$ & $p$ & $\widehat{B}$ & $p$ & $\widehat{B}$ & LCM & Fisher & Disc. & Binomial & CS1 & CS2B & Min & $T_\infty$ & $T_H$ & $T_H^{\mathrm{hybrid}}$ & Min \\
\midrule
RCT & 6795 & 141 & 185 & 0.0014 & .00 & .0026 & .00 & .0034 & .64 & 1.00 & .04 & .01 & .47 & .78 & .01 & .02 & .00 & .00$^{\dagger}$ & .00 \\
IV & 4962 & 279 & 262 & 0.0013 & .00 & .0033 & .00 & .0036 & 1.00 & 1.00 & .83 & .38 & .76 & .07 & .07 & .19 & .01 & .01$^{\dagger}$ & .01 \\
DID & 5316 & 237 & 293.5 & 0.0014 & .01 & .0016 & .02 & .0013 & .37 & 1.00 & .04 & .06 & .77 & .05 & .04 & .01 & .00 & .00$^{\dagger}$ & .00 \\
RDD & 2820 & 83 & $\cdot$ & $\cdot$ & .07 & .0000 & .06 & .0000 & .99 & 1.00 & .03 & .01 & .70 & .52 & .01 & .10 & .03 & .03$^{\dagger}$ & .03 \\
\bottomrule
\end{tabular}}

\medskip
\begin{minipage}{\linewidth}
Results for the Brodeur et al. (2020) data. We drop tests whose reported coefficient is missing; unlike Kudrin (2025), we keep those reported as zero.
De-rounding: none. Each $t$-statistic is the ratio of the reported coefficient and standard error, and reported $t$-statistics are used as printed.
The sample is the same for every de-rounding choice: we also drop $t$-statistics whose standard error is reported as zero or whose coefficient's rounding unit exceeds 0.5 times the smallest standard error consistent with the reported one.
Columns headed $J_0$ report $p$-values and breakdown statistics $\widehat{B}$ of the projection test with $J=40$; several values of $J_0$ are included to show low sensitivity to the smoothing parameter.
Med.\ $\nu$ is the median effective sample size as measured by the median number of clusters in a random sample of 100 articles for each causal inference method.
In ambiguous cases, the smallest cluster count was always used.
$\Delta_\nu=\frac{1}{100}\sum_{k=1}^{100}\sup_h\|\varphi(t-h)-t_{\nu_k,h}(t)\|_T$ is the average difference between the PDFs of the standard normal and Student-$t$ over the sample of 100 effective sample sizes collected (computed numerically).
Since the $p$-value for RDDs is insignificant, we did not compute $\Delta_\nu$ for them.
We form $\mathcal{X}$ using a grid of 3000 elements evenly spaced between $\pm 6.5$ plus the scalar 9999.
Discretization incurs error $\mathcal{E}(L,\delta)=0.00063$, which is already accounted for in $\widehat{B}$ and the $p$-value.
The discretization was chosen to be made of round numbers and result in error relatively small compared to $\Delta_\nu$.
EKW $p$-values are for the tests of Elliott et al. (2022): the least concave majorant test (LCM) and Fisher's test (Fisher) on $p$-values in $(0,0.15]$, the density discontinuity test at 0.05 (Disc.), the binomial test on $[0.04,0.05]$ (Binomial), and the Cox and Shi (2023) histogram-based tests on $(0,0.15]$ with 15 bins for monotonicity (CS1) and for 2-monotonicity with bounds (CS2B).
Kudrin (2025) $p$-values are for his KS-type test ($T_\infty$), histogram-based test ($T_H$) and hybrid histogram-based test ($T_H^{\mathrm{hybrid}}$), using the sinc kernel with $\kappa=3$ and, for $T_H$ and $T_H^{\mathrm{hybrid}}$, the points $0.01,0.02,\ldots,0.15$.
All three cluster by article: $T_H$ and $T_H^{\mathrm{hybrid}}$ use cluster-robust covariance matrices and $T_\infty$ uses an article-level bootstrap.
$^\dagger$The deconvolution estimate already satisfies the hybrid test's shape restrictions, so $T_H^{\mathrm{hybrid}}=T_H$.
Min is the smallest $p$-value in each group of tests.
The projection test and $T_\infty$ use 1000 bootstrap repetitions.
Computation used osqp package in R by Stellato et al. (2024). Data: Brodeur et al. (2020).
Transparency R package available at: \url{https://github.com/sfaridan/When_is_p_hacking_detectable}
\end{minipage}
\end{table}
